% Part: methods
% Chapter: proofs
% Section: inference-patterns

\documentclass[../../../include/open-logic-section]{subfiles}

\begin{document}

\olfileid[pt]{mth}{prf}{pat}

\olsection{Padrões de Inferência}

As provas são compostas de inferências individuais. Quando fazemos uma
inferência, normalmente a indicamos usando uma palavra como ``então,''
``assim,'' ou ``portanto.''  A inferência muitas vezes se apoia em um ou dois
fatos que já temos disponíveis em nossa prova---pode ser algo que
supusemos, ou algo que concluímos por uma inferência
anterior.  Para deixar claro, podemos rotular essas coisas, e na inferência
indicamos quais outras afirmações estamos usando na inferência.  Uma
inferência muitas vezes também conterá uma explicação de \emph{por que} nossa nova
conclusão decorre das coisas que a precedem.  Há
alguns padrões comuns de inferência que são usados com muita frequência em provas;
percorreremos alguns deles abaixo. Alguns padrões de inferência, como as provas
por indução, são mais elaborados (e serão discutidos mais adiante).

Já discutimos um padrão de inferência: desdobrar, ou
aplicar, uma definição. Quando desdobramos uma definição, apenas reenunciamos
algo que envolve o definiendum usando o definiens. Por
exemplo, suponha que já tenhamos estabelecido no curso de uma
prova que $D = E$ (a). Então podemos aplicar a definição de $=$ para conjuntos
e inferir: ``Assim, por definição, de (a), todo !!{element}
de~$D$ é !!a{element} de~$E$ e vice-versa.''

De modo um tanto confuso, muitas vezes não escrevemos a justificativa de uma
inferência quando de fato a fazemos, mas antes.  Suponha que ainda não
tenhamos provado que $D = E$, mas queremos.  Se $D = E$ é a
conclusão que buscamos, então podemos reenunciar esse objetivo também aplicando
a definição: para provar $D = E$ temos que provar que todo
!!{element} de~$D$ é !!a{element} de~$E$ e vice-versa.  Assim, nossa
prova terá a forma: (a) provar que todo !!{element} de~$D$ é
!!a{element} de~$E$; (b) todo !!{element} de~$E$ é !!a{element}
de~$D$; (c) portanto, de (a) e (b), pela definição de $=$, $D = E$.
Mas normalmente não escreveríamos desse jeito. Em vez disso, poderíamos escrever
algo como,
\begin{quote}
Queremos mostrar $D = E$. Pela definição de~$=$, isso equivale a mostrar
que todo !!{element} de~$D$ é !!a{element} de~$E$ e vice-versa.

(a) \dots (uma prova de que todo !!{element} de~$D$ é !!a{element}
de~$E$) \dots

(b) \dots (uma prova de que todo !!{element}
de~$E$ é !!a{element} de~$D$) \dots
\end{quote}

\subsection{Usando uma Conjunção}

Talvez o padrão de inferência mais simples seja o de tirar como
conclusão uma das conjuntivas de uma conjunção. Em outras palavras: se
supusemos ou já provamos que $p$ e~$q$, então temos o
direito de inferir que~$p$ (e também que~$q$).  Essa é uma inferência tão básica
que muitas vezes não é mencionada.  Por exemplo, uma vez que tenhamos
desdobrado a definição de $D = E$, estabelecemos que todo
!!{element} de~$D$ é !!a{element} de~$E$ e vice-versa. Disso
podemos concluir que todo !!{element} de~$E$ é !!a{element} de~$D$
(essa é a parte do ``vice-versa'').  

\subsection{Provando uma Conjunção}

Às vezes, o que lhe pedirão para provar terá a forma de uma
conjunção; pedirão a você que ``prove $p$ e $q$.'' Nesse caso,
você simplesmente tem que fazer duas coisas: provar $p$, e então provar $q$. Você
poderia dividir sua prova em duas seções e, para maior clareza, rotulá-las.
Quando estiver fazendo suas primeiras anotações, você pode escrever ``(1) Provar
$p$'' no topo da página, e ``(2) Provar $q$'' no meio da
página. (É claro que você pode não ser explicitamente solicitado a provar uma
conjunção, mas descobrir que sua prova exige que você prove uma
conjunção. Por exemplo, se lhe pedem para provar que $D = E$, você
descobrirá que, após desdobrar a definição de~$=$, você tem que
provar: todo !!{element} de~$D$ é !!a{element} de~$E$
\emph{e} todo !!{element} de~$E$ é !!a{element} de~$D$).

\subsection{Provando uma Disjunção}

Quando o que você está provando assume a forma de uma disjunção (isto é, é
uma afirmação da forma ``$p$ ou $q$''), basta mostrar que
um dos disjuntos é verdadeiro.  No entanto, basicamente nunca acontece
de qualquer um dos disjuntos simplesmente decorrer das suposições do seu
teorema. Mais frequentemente, as suposições do seu teorema são elas próprias
disjuntivas, ou você está mostrando que todas as coisas de certo tipo têm
uma de duas propriedades, mas algumas das coisas têm uma e outras
têm a outra propriedade.  É aqui que a prova por casos é
útil (veja abaixo).

\subsection{Prova Condicional}

Muitos teoremas que você encontrará estão na forma condicional (isto é, mostre
que se $p$ vale, então $q$ também é verdadeiro). Esses casos são bons e
fáceis de montar---basta supor o antecedente do condicional (neste
caso, $p$) e provar a conclusão~$q$ a partir dele.  Então, se o seu
teorema diz ``Se $p$ então $q$,'' você começa sua prova com ``suponha
$p$'' e, ao final, você deve ter provado~$q$.

Os condicionais podem ser enunciados de maneiras diferentes. Assim, em vez de ``Se $p$
então $q$,'' um teorema pode afirmar que ``$p$ somente se $q$,'' ``$q$ se
$p$,'' ou ``$q$, desde que $p$.'' Todos eles significam o mesmo e exigem
supor $p$ e provar~$q$ a partir dessa suposição.  Lembre-se de que um
bicondicional (``$p$ se e somente se (sse) $q$'') é, na verdade, dois
condicionais postos juntos: se $p$ então $q$, e se $q$ então~$p$. Tudo
o que você tem que fazer, então, são duas instâncias de prova condicional: uma para
o primeiro condicional e outra para o segundo. Às vezes,
no entanto, é possível provar uma afirmação ``sse'' encadeando
um monte de outras afirmações ``sse'', de modo que você comece com
``$p$'' e termine com ``$q$''---mas, nesse caso, você tem que garantir
que cada passo realmente seja um ``sse.''

\subsection{Afirmações Universais}

Usar uma afirmação universal é simples: se algo é verdadeiro para qualquer coisa,
é verdadeiro para cada coisa particular.  Então, se, digamos, a hipótese da
sua prova é $A \subseteq B$, isso significa (desdobrando a definição
de~$\subseteq$) que, para todo $x \in A$, $x \in B$. Assim, se você
já sabe que $z \in A$, você pode concluir $z \in B$.

Provar uma afirmação universal pode parecer um pouco complicado. Normalmente essas
afirmações assumem a seguinte forma: ``Se $x$ tem~$P$, então ele
tem~$Q$'' ou ``Todos os $P$ são $Q$.'' É claro que pode não se ajustar a essa
forma perfeitamente, e é preciso um pouco de prática para descobrir o que
exatamente lhe pedem para provar. Mas: muitas vezes temos que provar que todos os objetos
com alguma propriedade têm uma certa outra propriedade.

A maneira de provar uma afirmação universal é introduzir nomes ou variáveis
para as coisas que têm uma propriedade e então mostrar que elas também
têm a outra propriedade.  Poderíamos expressar isso dizendo que, para provar
algo para \emph{todos} os $P$, você tem que prová-lo para um
$P$~\emph{arbitrário}. E o nome introduzido é um nome para um
$P$~arbitrário.  Normalmente usamos letras isoladas como esses nomes para
coisas arbitrárias, e as letras costumam seguir convenções: por exemplo,
usamos $n$ para números naturais, $!A$ para !!{formula}s, $A$ para conjuntos, $f$
para funções, etc.

O truque é manter a generalidade ao longo da prova. Você começa
supondo que um objeto arbitrário (``$x$'') tem a propriedade~$P$, e
mostra (com base apenas em definições ou no que lhe é permitido supor)
que $x$ tem a propriedade~$Q$. Como você não estipulou o que
$x$ é especificamente, a não ser que ele tem a propriedade $P$, então você pode
afirmar que tudo que tem $P$ tem a propriedade~$Q$. Em suma, $x$ é um
substituto para \emph{todas} as coisas com a propriedade~$P$.

\begin{prop}
  Para todos os conjuntos $A$ e $B$, $A \subseteq A \cup B$.
\end{prop}

\begin{proof}
  Sejam $A$ e $B$ conjuntos arbitrários.  Queremos mostrar que $A
  \subseteq A \cup B$. Pela definição de $\subseteq$, isso equivale a:
  para todo $x$, se $x \in A$ então $x \in A \cup B$. Então seja $x \in A$
  !!a{element} arbitrário de~$A$. Temos que mostrar que $x \in A
  \cup B$. Como $x \in A$, $x \in A$ ou $x \in B$. Assim, $x \in
  \Setabs{x}{x \in A \lor x \in B}$. Mas isso, pela definição de $\cup
  $, significa $x \in A \cup B$.
\end{proof}


\subsection{Prova por Casos}

Suponha que você tenha uma disjunção como suposição ou como conclusão já
estabelecida---você supôs ou provou que $p$ ou $q$ é
verdadeira.  Você quer provar $r$.  Você faz isso em dois passos: primeiro você
supõe que $p$ é verdadeira, e prova~$r$, e depois supõe que $q$ é
verdadeira e prova~$r$ novamente.  Isso funciona porque supomos ou sabemos que
uma das duas alternativas vale. Os dois passos estabelecem que qualquer
uma delas é suficiente para a verdade de~$r$.  (Se ambas são verdadeiras, temos
não uma, mas duas razões para que $r$~seja verdadeira. Não é necessário
provar separadamente que $r$~é verdadeira supondo tanto $p$ quanto~$q$.)  Para
indicar o que estamos fazendo, anunciamos que ``distinguimos casos.''
Por exemplo, suponha que saibamos que $x \in B \cup C$.  $B \cup C$ é
definido como $\Setabs{x}{x \in B \text{ ou } x \in C}$. Em outras palavras,
por definição, $x \in B$ ou $x \in C$. Provaríamos que $x \in A$
a partir disso supondo primeiro que $x \in B$, e provando $x \in A$ a partir
dessa suposição, e então supondo $x \in C$, e novamente provando $x \in A$
a partir disso.  Você escreveria ``Distinguimos casos'' sob a
suposição, depois ``Caso (1): $x \in B$'' abaixo, e ``Caso (2):
$x \in C$'' na metade da página. Então você prosseguiria preenchendo a metade
superior e a metade inferior da página.

A prova por casos é especialmente útil se o que você está provando é, ele mesmo,
disjuntivo. Eis um exemplo simples:

\begin{prop}
Suponha $B \subseteq D$ e $C \subseteq E$. Então $B \cup C \subseteq
D \cup E$.
\end{prop}

\begin{proof}
  Suponha (a) que $B \subseteq D$ e (b) $C \subseteq E$. Por
  definição, qualquer $x \in B$ também está $\in D$ (c) e qualquer $x \in C$ também
  está $\in E$ (d).  Para mostrar que $B \cup C \subseteq D \cup E$,
  temos que mostrar que se $x \in B \cup C$ então $x \in D \cup E$ (pela
  definição de $\subseteq$). $x \in B \cup C$ sse $x \in B$ ou $x \in
  C$ (pela definição de~$\cup$). De modo semelhante, $x \in D \cup E$ sse $x \in
  D$ ou $x \in E$. Então, temos que mostrar: para qualquer $x$, se $x \in B$ ou
  $x \in C$, então $x \in D$ ou $x \in E$.

  \begin{quote}
  Até aqui só desdobramos definições!{} Reformulamos nossa
  proposição sem $\subseteq$ e $\cup$ e ficamos tentando
  provar uma afirmação condicional universal. Pelo que discutimos
  acima, isso é feito supondo que $x$ é algo a respeito do qual
  supomos que a parte do ``se'' é verdadeira, e passaremos a mostrar que a
  parte do ``então'' também é verdadeira. Em outras palavras, suporemos que $x
  \in B$ ou $x \in C$ e mostraremos que $x \in D$ ou $x \in
  E$.\footnote{Este parágrafo apenas explica o que estamos fazendo---não
    é parte da prova, e você não precisa entrar em todo esse
    detalhe quando escreve suas próprias provas.}
  \end{quote}

  Suponha que $x \in B$ ou $x \in C$. Temos que mostrar que $x \in D$
  ou $x \in E$. Distinguimos casos.

  Caso 1: $x \in B$. Por (c), $x \in D$. Assim, $x \in D$ ou $x \in
  E$. (Aqui fizemos a inferência discutida na subseção
  anterior!)
      
  Caso 2: $x \in C$. Por (d), $x \in E$. Assim, $x \in D$ ou $x \in E$.
 \end{proof}


\subsection{Provando uma Afirmação de Existência}

Quando lhe pedem para provar uma afirmação de existência, a questão normalmente será
da forma ``prove que há um~$x$ tal que $\dots x \dots$'',
isto é, que algum objeto tem a propriedade descrita por ``$\dots x
\dots$''. Nesse caso, você terá que identificar um objeto adequado e mostrar
que ele tem a propriedade exigida.  Isso parece direto, mas uma
prova desse tipo pode ser complicada. Tipicamente, ela envolve
\emph{construir} ou \emph{definir} um objeto e provar que o
objeto assim definido tem a propriedade exigida. Encontrar o objeto certo
pode ser difícil, provar que ele tem a propriedade exigida pode ser difícil,
e às vezes é até complicado mostrar que você teve sucesso em
definir um objeto!{}

Em geral, você escreveria isso especificando o objeto, por exemplo, ``seja
$x$ \dots'' (em que \dots{} especifica qual objeto você tem em
mente), possivelmente provando que $\dots$ de fato descreve um objeto que
existe, e então passaria a mostrar que $x$ tem a propriedade~$Q$. Eis um
exemplo simples.

\begin{prop}
  Suponha que $x \in B$. Então há um~$A$ tal que $A \subseteq
  B$ e $A \neq \emptyset$.
\end{prop}

\begin{proof}
  Suponha $x \in B$. Seja $A = \{x\}$.
  \begin{quote}
    Aqui definimos o conjunto~$A$ enumerando seus
    !!{element}s. Como supomos que $x$ é um objeto, e sempre podemos
    formar um conjunto enumerando seus !!{element}s, não temos
    que mostrar que tivemos sucesso em definir um conjunto~$A$ aqui.  No entanto,
    ainda temos que mostrar que $A$ tem as propriedades exigidas pela
    proposição. A prova não está completa sem isso!{}
  \end{quote}
  Como $x \in A$, $A \neq \emptyset$.
  \begin{quote}
    Isso se apoia na definição de $A$ como $\{x\}$ e nos fatos óbvios
    de que $x \in \{x\}$ e $x \notin \emptyset$.
  \end{quote}
  Como $x$ é o único !!{element} de~$\{x\}$, e $x \in B$, todo
  !!{element} de~$A$ é também !!a{element} de~$B$. Pela definição
  de~$\subseteq$, $A \subseteq B$.
\end{proof}

\subsection{Usando Afirmações de Existência}

Suponha que você saiba que alguma afirmação de existência é verdadeira (você a provou,
ou ela é uma hipótese que você pode usar), digamos, ``para algum~$x$, $x \in A$'' ou
``há um $x \in A$.''  Se você quiser usá-la em sua prova, você
pode simplesmente fingir que tem um nome para uma das coisas que sua
hipótese diz existir. Como $A$ contém pelo menos uma coisa, há
coisas a que esse nome poderia se referir. É claro que você pode não ser
capaz de escolher uma ou de descrevê-la mais (a não ser que ela é
$\in A$). Mas, para os fins da prova, você pode fingir que
a escolheu e dar-lhe um nome. É importante escolher um
nome que você ainda não tenha usado (ou que apareça nas suas
hipóteses), caso contrário as coisas podem dar errado. Em sua prova, você
indica isso passando de ``para algum $x$, $x \in A$'' para ``Seja $a
\in A$.''  Agora você pode raciocinar sobre~$a$, usar algumas outras hipóteses, etc.,
até chegar a uma conclusão, $p$. Se $p$ não menciona mais~$a$, $p$ é
independente da suposição de que $a \in A$, e você mostrou que ele
decorre apenas da suposição ``para algum $x$, $x \in A$.''

\begin{prop}
Se $A \neq \emptyset$, então $A \cup B \neq \emptyset$.
\end{prop}

\begin{proof}
  Suponha $A \neq \emptyset$. Então, para algum $x$, $x \in A$. 
  \begin{quote}
    Aqui apenas reenunciamos a hipótese da
    proposição. Essa hipótese, isto é, $A \neq \emptyset$, esconde uma
    afirmação existencial, à qual você só chega desdobrando algumas
    definições. A definição de $=$ nos diz que $A = \emptyset$
    sse todo $x \in A$ também está $\in \emptyset$ e todo $x \in
    \emptyset$ também está $\in A$. Negando ambos os lados, obtemos: $A \neq
    \emptyset$ sse ou algum $x \in A$ está $\notin \emptyset$ ou algum
    $x \in \emptyset$ está $\notin A$. Como nada está $\in \emptyset$,
    o segundo disjunto nunca pode ser verdadeiro, e ``$x \in A$ e $x
    \notin \emptyset$'' se reduz a apenas $x \in A$. Então $A \neq
    \emptyset$ sse, para algum $x$, $x \in A$. Essa é uma afirmação de
    existência. Agora usamos essa afirmação de existência introduzindo um nome para
    um dos !!{element}s de~$A$:
  \end{quote}
  Seja $a \in A$.
  \begin{quote}
    Agora introduzimos um nome para uma das coisas~$\in A$. Continuaremos
    a argumentar sobre~$a$, mas teremos o cuidado de supor apenas
    que $a \in A$ e nada mais:
  \end{quote}
  Como $a \in A$, $a \in A \cup B$, pela definição de~$\cup$. Então, para
  algum $x$, $x \in A \cup B$, isto é, $A \cup B \neq \emptyset$.
  \begin{quote}
    Naquele último passo, fomos de ``$a \in A \cup B$'' para ``para algum
    $x$, $x \in A \cup B$.'' Isso não menciona mais $a$, então
    sabemos que ``para algum $x$, $x \in A \cup B$'' decorre de ``para
    algum $x$, $x \in A$'' apenas. Mas isso significa que $A \cup B \neq
    \emptyset$.
  \end{quote}
\end{proof}

Talvez seja boa prática manter variáveis ligadas como ``$x$'' separadas
de nomes hipotéticos como $a$, como fizemos. Na prática, porém,
muitas vezes não o fazemos e apenas usamos $x$, assim:
\begin{quote}
Suponha $A \neq \emptyset$, isto é, há um $x \in A$. Pela definição
de $\cup$, $x \in A \cup B$. Então $A \cup B \neq \emptyset$.
\end{quote}
No entanto, quando você faz isso, tem que ter um cuidado extra para usar
$x$ e $y$ diferentes para afirmações existenciais diferentes. Por
exemplo, a seguinte \emph{não} é uma prova correta de ``Se $A \neq
\emptyset$ e $B \neq \emptyset$ então $A \cap B \neq \emptyset$''
(o que não é verdade).
\begin{quote}
Suponha $A \neq \emptyset$ e $B \neq \emptyset$. Então, para algum $x$, $x
\in A$ e também, para algum $x$, $x \in B$. Como $x \in A$ e $x \in
B$, $x \in A \cap B$, pela definição de~$\cap$. Então $A \cap B \neq
\emptyset$.
\end{quote}
Você consegue identificar onde ocorre o passo incorreto e explicar por que o
resultado não vale?

\end{document}
